\chapter{Заряженные частицы, преобразования полей и релятивистские волны}

В прошлый раз мы начали выяснять, как движется заряженная частица в электрическом и магнитном полях, когда учитываются принципы теории относительности. Уже знакомый закон силы Лоренца содержит подсказку: его магнитный член зависит от скорости частицы, поэтому разные наблюдатели не могут по отдельности прийти к согласию о том, что является электрическим, а что магнитным. Наша задача — сформулировать этот закон так, чтобы каждый инерциальный наблюдатель пользовался одним и тем же уравнением, хотя поля и скорость частицы меняются при переходе от одной системы отсчета к другой.

Если не оговорено иное, мы используем единицы, в которых \(c=1\), а сигнатура метрики имеет вид \((-,-,-,+)\). Компоненты пространства-времени упорядочены так, что сначала идут три пространственные компоненты:
\begin{equation}
  \mu=(1,2,3,0)=(x,y,z,t).
\end{equation}
Латинские индексы \(m,n\) пробегают только три пространственных направления.

\section{Что остается от ньютоновского закона силы}

В ньютоновской механике
\begin{equation}
  \boldsymbol F=m\boldsymbol a
\end{equation}
имеет один и тот же вид во всех инерциальных системах отсчета, потому что галилеевы наблюдатели согласны относительно ускорения. Они также приписывают частице одну и ту же массу и одну и ту же силу. Как только кинематика меняется, нет оснований ожидать, что именно это соотношение между силой и обычным ускорением сохранится без изменений.

Электромагнетизм и гравитация уже предупреждают нас, что релятивистские обобщения не обязаны быть похожими друг на друга. Закон Кулона и ньютоновский закон тяготения оба являются законами обратных квадратов, но гравитацию нельзя естественным образом включить в специальную теорию относительности без более глубокого изменения геометрии пространства-времени. С электромагнетизмом это возможно. Более того, магнетизм сам по себе является релятивистским эффектом: когда множители \(c\) восстановлены, магнитная сила содержит множитель скорости, деленной на \(c\), и исчезает в формальном пределе \(c\to\infty\).

То, что в ньютоновском законе сохраняется, — это утверждение, что сила есть скорость изменения импульса. При постоянной массе покоя
\begin{equation}
  \boldsymbol F
  =m\frac{d\boldsymbol v}{dt}
  =\frac{d}{dt}(m\boldsymbol v)
  =\frac{d\boldsymbol p}{dt}.
\end{equation}
Теория относительности изменит связь между импульсом и скоростью, но сохранит этот способ организации уравнения.

\begin{figure}[H]
\centering
\includegraphics[width=0.64\textwidth]{lecture_07_figure_01.png}
\caption{Сила переписывается в виде \(d\boldsymbol p/dt\), в то время как на доске записывается электрическая часть закона Лоренца.}
\lectureframe{00:08:20}
\end{figure}

В единицах \(c=1\) дорелятивистский закон силы Лоренца имеет вид
\begin{equation}
  \frac{d\boldsymbol p}{dt}
  =q\bigl(\boldsymbol E+\boldsymbol v\times\boldsymbol B\bigr).
  \label{eq:lorentz-three}
\end{equation}
Оба поля могут зависеть от положения и времени. Нерелятивистский предел теперь есть предел \(v\ll 1\), а не предел, в котором числовая константа \(c=1\) каким-то образом меняется. Электрический член в \eqref{eq:lorentz-three} уже является обычным вектором. Чтобы понять магнитный член в форме, которую можно перенести в пространство-время, надо снова посмотреть на векторное произведение.

\section{Векторное произведение как антисимметричный тензор}

Возьмем два трехмерных вектора \(\boldsymbol A\) и \(\boldsymbol B\). Их внешнее произведение
\begin{equation}
  T_{mn}=A_mB_n
\end{equation}
является тензором ранга два с девятью компонентами. Вообще говоря, никакой симметрией он не обладает. Антисимметризация дает
\begin{equation}
  K_{mn}=A_mB_n-A_nB_m,
  \qquad K_{mn}=-K_{nm}.
  \label{eq:antisym-product}
\end{equation}
Его диагональные элементы равны нулю, а каждый элемент ниже диагонали определяется соответствующим элементом выше нее. Следовательно, антисимметричный тензор \(3\times3\) имеет только три независимые компоненты.

Этот подсчет наводит на мысль, но сам по себе недостаточен. Три величины не автоматически преобразуются как вектор. Однако в трех измерениях три независимых элемента \eqref{eq:antisym-product} можно отождествить с компонентами векторного произведения \(\boldsymbol C=\boldsymbol A\times\boldsymbol B\). Сначала на доске \(C_2\) временно помещается в правый верхний угол:

\begin{figure}[H]
\centering
\includegraphics[width=0.75\textwidth]{lecture_07_figure_02.png}
\caption{Антисимметричное произведение в тот момент, когда его три независимых элемента сопоставляются с компонентами \(\boldsymbol C\). Обозначение в правом верхнем углу здесь еще предварительное.}
\lectureframe{00:17:33}
\end{figure}

Циклическая перестановка \(1\to2\to3\to1\) дает формулы для компонент
\begin{align}
  C_3 &= A_1B_2-A_2B_1, \\
  C_1 &= A_2B_3-A_3B_2, \\
  C_2 &= A_3B_1-A_1B_3.
  \label{eq:cross-components}
\end{align}
Последняя строка показывает, что элемент в правом верхнем углу равен \(-C_2\), а не \(C_2\). Эта поправка составляет часть рассуждения, а не деталь, которую следует скрыть.

\begin{figure}[H]
\centering
\includegraphics[width=0.62\textwidth]{lecture_07_figure_03.png}
\caption{Завершенные циклические тождества для компонент, которые исправляют временный знак в антисимметричном массиве.}
\lectureframe{00:21:54}
\end{figure}

Получающаяся матрица имеет вид
\begin{equation}
  K_{mn}=
  \begin{pmatrix}
    0 & C_3 & -C_2 \\
    -C_3 & 0 & C_1 \\
    C_2 & -C_1 & 0
  \end{pmatrix}.
  \label{eq:cross-matrix}
\end{equation}

За этим обозначением стоит полезная геометрия. Компонента с индексами \(1,2\) естественным образом относится к плоскости \(xy\) и отождествляется с векторной компонентой вдоль перпендикулярной оси \(z\). Один и тот же ориентированный объект можно описывать либо плоскостью, либо стрелкой, нормальной к этой плоскости. Описание через стрелку специфично для трех измерений; в более высоких размерностях антисимметричный тензор ранга два по-прежнему существует, но число его компонент уже не совпадает с числом компонент вектора. Поэтому обобщается именно тензор, а не перпендикулярная стрелка.

Эквивалентно, после выбора ориентации,
\begin{equation}
  C_k=\frac12\epsilon_{kmn}K_{mn}.
\end{equation}
Необходимость выбрать ориентацию лежит в основе правила правой руки и привычного различия между полярными и аксиальными векторами.
\editorialnote{Формула Леви--Чивиты и терминология аксиальных векторов — это стандартные обозначения для геометрии, в которой сопоставляются плоскость и нормаль, явно развернутой в лекции.}

\section{Магнитное поле в тензорной форме}

То же двойственное описание применимо и к магнитному полю. Его антисимметричный пространственный тензор определим как
\begin{equation}
  B_{mn}=
  \begin{pmatrix}
    0 & -B_3 & B_2 \\
    B_3 & 0 & -B_1 \\
    -B_2 & B_1 & 0
  \end{pmatrix}.
  \label{eq:magnetic-tensor}
\end{equation}
Дополнительный общий знак по сравнению с \eqref{eq:cross-matrix} является соглашением. После того как он выбран, им нужно пользоваться последовательно; изменение знака у всех трех векторных компонент описывало бы то же самое отождествление тензора с вектором, но при противоположном соглашении об ориентации.

\begin{figure}[H]
\centering
\includegraphics[width=0.72\textwidth]{lecture_07_figure_04.png}
\caption{Магнитное поле, записанное и как антисимметричный пространственный массив, и через его три векторные компоненты.}
\lectureframe{00:24:40}
\end{figure}

Теперь возьмем третью компоненту векторного произведения:
\begin{align}
  (\boldsymbol v\times\boldsymbol B)_3
  &=v_1B_2-v_2B_1 \\
  &=v_1B_{13}+v_2B_{23}+v_3B_{33}.
\end{align}
Последний член можно добавить, потому что \(B_{33}=0\). Циклическая перестановка свободного индекса дает компактное правило
\begin{equation}
  (\boldsymbol v\times\boldsymbol B)_m=v_nB_{nm},
  \label{eq:cross-contraction}
\end{equation}
где по повторяющемуся индексу \(n\) ведется суммирование. Следовательно, магнитная сила равна
\begin{equation}
  F_m^{(B)}=qv_nB_{nm}.
\end{equation}

\begin{figure}[H]
\centering
\includegraphics[width=0.78\textwidth]{lecture_07_figure_05.png}
\caption{Третья компонента \(\boldsymbol v\times\boldsymbol B\) и ее общее выражение как свертка с \(B_{nm}\).}
\lectureframe{00:28:30}
\end{figure}

\begin{classroomqa}
\audiencequestion{Обозначает ли свободный индекс \(m\) отдельную магнитную компоненту силы?}
\lecturerresponse{Нет. Он обозначает ту пространственную компоненту силы, которая в данный момент вычисляется. Для целей настоящего рассуждения электрическое поле просто положено равным нулю, и каждый индекс здесь пробегает только три пространственных направления.}
\end{classroomqa}

Тензорная запись делает свойство магнитного поля не совершать работу почти автоматическим. Продифференцируем квадрат модуля скорости:
\begin{align}
  \frac{d}{dt}(v_mv_m)
  &=2\dot v_m v_m \\
  &\propto 2(\boldsymbol v\times\boldsymbol B)_m v_m \\
  &=2v_nB_{nm}v_m.
  \label{eq:speed-derivative}
\end{align}
Члены с \(m=n\) исчезают, потому что диагональ \(B_{nm}\) равна нулю. У каждого внедиагонального члена есть парный член с переставленными \(m\) и \(n\); эти два сокращаются, поскольку \(v_nv_m\) симметричен, а \(B_{nm}\) антисимметричен. Следовательно,
\begin{equation}
  \frac{d}{dt}(v_mv_m)=0.
\end{equation}
Чисто магнитная сила искривляет траекторию, не изменяя ни модуля скорости, ни кинетической энергии.

\begin{figure}[H]
\centering
\includegraphics[width=0.68\textwidth]{lecture_07_figure_06.png}
\caption{Парные члены в \(v_nB_{nm}v_m\) сокращаются, потому что магнитный тензор меняет знак при перестановке своих индексов.}
\lectureframe{00:35:40}
\end{figure}

Это доказательство намеренно более развернуто, чем просто сказать, что векторное произведение перпендикулярно своему первому множителю. Последнее утверждение специфично для трехмерного векторного представления. Свертка симметричной пары с антисимметричным тензором — вот тот аргумент, который сохранится и в четырех измерениях.

\begin{classroomqa}
\audiencequestion{Должно ли магнитное поле не зависеть от времени, чтобы модуль скорости оставался постоянным?}
\lecturerresponse{В приведенной выше алгебре поле нигде не дифференцируется и не требуется \(\partial_t\boldsymbol B=0\). Сам по себе магнитный член все равно не совершает работы. Однако в электродинамике изменяющееся магнитное поле порождает электрическое поле, и это электрическое поле может изменить скорость частицы.}
\end{classroomqa}

\section{Почему электрическое и магнитное поля должны смешиваться}

Смешивание \(\boldsymbol E\) и \(\boldsymbol B\) можно увидеть еще до всяких четырехмерных вычислений. Рассмотрим медленное движение, так что обычных ньютоновских утверждений об ускорении достаточно. Пусть магнитное поле направлено вдоль оси \(y\), а заряженная частица движется вдоль оси \(x\). Тогда магнитная сила направлена вдоль оси \(z\):
\begin{equation}
  \boldsymbol F=q\boldsymbol v\times\boldsymbol B.
\end{equation}

\begin{figure}[H]
\centering
\includegraphics[width=0.62\textwidth]{lecture_07_figure_07.png}
\caption{Набросок на доске для частицы, движущейся поперек магнитного поля; три направления используются для сравнения описаний в лабораторной и сопутствующей системах отсчета.}
\lectureframe{00:42:20}
\end{figure}

Эту геометрию можно перерисовать без перспективы, свойственной рисунку на доске:
\begin{figure}[H]
\centering
\begin{tikzpicture}[scale=1.05,>=Stealth]
  \draw[->] (0,0) -- (3.8,0) node[right] {\(x\)};
  \draw[->] (0,0) -- (-1.9,2.4) node[above left] {\(y\)};
  \draw[->] (0,0) -- (0,3.1) node[above] {\(z\)};
  \draw[->,very thick] (0,0) -- (2.7,0)
    node[midway,below] {\(\boldsymbol v\)};
  \draw[->,very thick] (0,0) -- (-1.35,1.7)
    node[midway,above left] {\(\boldsymbol B\)};
  \draw[->,very thick] (0,0) -- (0,2.35)
    node[midway,right] {\(\boldsymbol F,\ \boldsymbol E'\)};
\end{tikzpicture}
\caption{В лабораторной системе отсчета сила, направленная вдоль \(z\), является магнитной; в мгновенно сопутствующей системе отсчета та же самая сила должна быть электрической.}
\end{figure}

Теперь введём наблюдателя, движущегося вместе с частицей в этот момент. Скорость частицы в новой системе отсчёта равна нулю, поэтому магнитный член \(q\boldsymbol v'\times\boldsymbol B'\) исчезает. Однако в пределе малых скоростей оба наблюдателя должны согласиться, что частица имеет \(z\)-ускорение. Если закон силы Лоренца имеет один и тот же вид в обеих системах, то сопутствующий наблюдатель должен приписать силу электрическому полю \(\boldsymbol E'\) вдоль \(z\). Следовательно, поле, которое в лаборатории было чисто магнитным, в движущейся системе не является чисто магнитным.

Именно по этой физической причине стоит искать единый объект, компоненты которого включают оба поля. При переходе к другой системе отсчёта происходит не просто поворот \(\boldsymbol E\) среди его собственных трёх компонент и \(\boldsymbol B\) среди его собственных трёх компонент; при этом два набора перемешиваются.

\section{Тензор электромагнитного поля}

Вместе электрическое и магнитное поля имеют 6 компонент. Четырёхмерный антисимметричный тензор тоже имеет 6 независимых компонент, поскольку
\begin{equation}
  \frac{4(4-1)}{2}=6.
\end{equation}
Это не означает, что антисимметричный тензор — буквально единственный набор из 6 величин, который можно было бы придумать: 6 скаляров или вектор и 2 скаляра тоже содержали бы 6 чисел. Это естественный геометрический объект, компоненты которого могут смешиваться требуемым образом.

Чтобы явно сохранить принятую на доске конвенцию, выпишем контравариантный тензор в порядке, где сначала идут пространственные индексы, \((1,2,3,0)\):
\begin{equation}
  F^{\mu\nu}=
  \begin{pmatrix}
    0 & -B_3 & B_2 & E_1 \\
    B_3 & 0 & -B_1 & E_2 \\
    -B_2 & B_1 & 0 & E_3 \\
    -E_1 & -E_2 & -E_3 & 0
  \end{pmatrix}.
  \label{eq:l07-field-tensor}
\end{equation}
Левый верхний блок \(3\times3\) — это в точности магнитный тензор \eqref{eq:magnetic-tensor}; компоненты с одним пространственным и одним временным индексом содержат электрическое поле. Остальные элементы задаются антисимметрией.

\begin{figure}[H]
\centering
\includegraphics[width=0.72\textwidth]{lecture_07_figure_08.png}
\caption{Магнитный пространственный блок и электрический временной-пространственный столбец, собранные в четырёхмерный тензор поля.}
\lectureframe{00:46:35}
\end{figure}

\begin{classroomqa}
\audiencequestion{Когда у этих компонент поля должны стоять верхние, а не нижние индексы?}
\lecturerresponse{Если речь идёт только о трёх пространственных обозначениях, на доске это различие иногда опускается при отождествлении привычных полей Maxwell--Faraday. В ковариантном вычислении различие необходимо восстановить: поднятие или опускание пространственного индекса меняет знак в выбранной метрической конвенции.}
\end{classroomqa}

Инвариантным является не конкретное расположение знаков на доске, а антисимметричный пространственно-временной тензор вместе с согласованной метрической и индексной конвенцией. Обозначение \(F\) может наводить на мысль о Faraday, хотя лекция не берётся с уверенностью утверждать такое историческое происхождение. Einstein вывел преобразование электрических и магнитных компонент в 1905; компактный четырёхмерный тензорный язык появился в более поздней пространственно-временной формулировке, связанной с Minkowski.

\section{Четырёхскорость и четырёхсила}

Сторону уравнения, относящуюся к частице, тоже нужно сделать ковариантной. Собственное время \(\tau\) инвариантно, поэтому дифференцирование четырёхвектора положения по нему даёт ещё один четырёхвектор:
\begin{equation}
  u^\mu=\frac{dx^\mu}{d\tau}
  =\gamma(v_x,v_y,v_z,1),
  \qquad
  \gamma=\frac{1}{\sqrt{1-v^2}}.
  \label{eq:four-velocity}
\end{equation}
Последняя компонента следует из \(dt/d\tau=\gamma\). Четырёхимпульс равен
\begin{equation}
  p^\mu=mu^\mu,
\end{equation}
его пространственная часть — это релятивистский импульс \(\boldsymbol p=m\gamma\boldsymbol v\).

\begin{figure}[H]
\centering
\includegraphics[width=0.75\textwidth]{lecture_07_figure_09.png}
\caption{Четырёхскорость, её компоненты, фактор Лоренца и четырёхимпульс.}
\lectureframe{00:50:30}
\end{figure}

Поскольку \(p^\mu\) — четырёхвектор, а \(\tau\) — скаляр, естественная релятивистская левая часть имеет вид
\begin{equation}
  f^\mu\equiv\frac{dp^\mu}{d\tau}.
  \label{eq:four-force-def}
\end{equation}
Четырёхсилу нельзя выбирать произвольно. У четырёхскорости фиксированная норма,
\begin{equation}
  u_\mu u^\mu=1.
\end{equation}
Дифференцируя вдоль мировой линии, получаем
\begin{align}
  0
  &=\frac{d}{d\tau}(u_\mu u^\mu) \\
  &=\frac{du_\mu}{d\tau}u^\mu
    +u_\mu\frac{du^\mu}{d\tau} \\
  &=2u_\mu\frac{du^\mu}{d\tau}.
\end{align}
При постоянной массе отсюда следует
\begin{equation}
  f^\mu u_\mu=0.
  \label{eq:force-orthogonal}
\end{equation}
Таким образом, четырёхсила ортогональна четырёхскорости в метрике Минковского. Это чисто кинематическое ограничение, вытекающее из нормировки \(u^\mu\), а не особое свойство электромагнетизма.

\begin{figure}[H]
\centering
\includegraphics[width=0.64\textwidth]{lecture_07_figure_10.png}
\caption{Единичная норма четырёхскорости и вытекающая из неё ортогональность четырёхсилы и четырёхскорости.}
\lectureframe{00:56:50}
\end{figure}

Трёхмерное магнитное вычисление уже подсказало приём, позволяющий удовлетворить \eqref{eq:force-orthogonal}. Свернём антисимметричный тензор с одним множителем скорости:
\begin{equation}
  \boxed{
  \frac{dp^\mu}{d\tau}=qF^{\mu\nu}u_\nu.}
  \label{eq:covariant-lorentz}
\end{equation}
Действительно,
\begin{equation}
  f^\mu u_\mu
  =qF^{\mu\nu}u_\nu u_\mu=0,
\end{equation}
поскольку \(F^{\mu\nu}\) антисимметричен, а \(u_\nu u_\mu\) симметричен.
Это не единственный мыслимый релятивистский закон силы, но это самый простой закон, линейный и по электромагнитному полю, и по скорости.

\begin{figure}[H]
\centering
\includegraphics[width=0.72\textwidth]{lecture_07_figure_11.png}
\caption{Ковариантное уравнение силы Лоренца слева и его обращающаяся в нуль свёртка со второй четырёхскоростью справа.}
\lectureframe{01:00:00}
\end{figure}

\section{Восстановление обычной силы Лоренца}

Уравнение, выглядящее ковариантно, всё равно должно воспроизводить известный закон силы. При метрической конвенции, заданной в начале,
\begin{equation}
  u_\mu=\gamma(-v_x,-v_y,-v_z,1).
\end{equation}
Для \(x\)-компоненты из \eqref{eq:l07-field-tensor} и \eqref{eq:covariant-lorentz} получается
\begin{align}
  \frac{dp^1}{d\tau}
  &=q\bigl(F^{12}u_2+F^{13}u_3+F^{10}u_0\bigr) \\
  &=q\gamma\bigl(B_3v_y-B_2v_z+E_1\bigr) \\
  &=q\gamma\bigl[\boldsymbol E+\boldsymbol v\times
    \boldsymbol B\bigr]_x.
  \label{eq:x-force}
\end{align}
Остальные пространственные компоненты получаются циклической перестановкой индексов.

\begin{classroomqa}
\audiencequestion{Где в разложении по компонентам находится член \(qE_x\)?}
\lecturerresponse{Это член \(qF^{10}u_0\). Компонента тензора поля с одним пространственным и одним временным индексом соответствует компоненте электрического поля, и \(u_0=\gamma\).}
\end{classroomqa}

Множитель \(\gamma\) справа — не дополнительная сила. Поскольку
\begin{equation}
  \frac{dp^i}{d\tau}=\frac{dt}{d\tau}\frac{dp^i}{dt}
  =\gamma\frac{dp^i}{dt},
\end{equation}
он сокращается в обеих частях \eqref{eq:x-force}. Поэтому точное пространственное уравнение снова имеет вид
\begin{equation}
  \frac{d\boldsymbol p}{dt}
  =q\bigl(\boldsymbol E+\boldsymbol v\times\boldsymbol B\bigr),
  \qquad
  \boldsymbol p=m\gamma\boldsymbol v.
  \label{eq:exact-three-force}
\end{equation}
Привычная форма сохраняется, но релятивистский характер скрыт в импульсе.

\begin{figure}[H]
\centering
\includegraphics[width=0.75\textwidth]{lecture_07_figure_12.png}
\caption{Трёхмерный закон Лоренца, восстановленный после удаления общего множителя собственного времени; импульс сохраняет свой множитель \(\gamma\).}
\lectureframe{01:06:30}
\end{figure}

Временная компонента несёт соответствующее утверждение об энергии,
\begin{equation}
  \frac{dp^0}{dt}=q\boldsymbol E\cdot\boldsymbol v,
\end{equation}
так что магнитное поле работы не совершает, тогда как электрическое поле может изменять энергию частицы.
\editorialnote{Проверка временной компоненты — прямое дополнение к пространственному вычислению в лекции; она включена, чтобы показать то же свойство отсутствия работы в четырёхмерных обозначениях.}

\begin{classroomqa}
\audiencequestion{Этот вывод — математика или физика?}
\lecturerresponse{И то и другое. Тензорная алгебра даёт компактную форму, но существенное физическое содержание вошло через принцип, согласно которому все инерциальные системы отсчёта подчиняются одним и тем же законам, и через предположение, что уравнения движения имеют второй порядок. Закон второго порядка предсказывает будущее по положению и скорости; закон третьего порядка потребовал бы включить ускорение в число начальных данных. Каким типом закона пользуется природа, говорит нам эксперимент, а не обозначения.}
\end{classroomqa}

\section{Как преобразуются поля}

Уравнение \eqref{eq:covariant-lorentz} имеет один и тот же вид в любой системе отсчёта, потому что обе его части — четырёхвекторы. Сам тензор поля преобразуется с одной матрицей Лоренца на каждый индекс:
\begin{equation}
  F'^{\mu\nu}
  =\Lambda^\mu{}_{\rho}\Lambda^\nu{}_{\sigma}F^{\rho\sigma}.
  \label{eq:tensor-transform}
\end{equation}
На доске это правило выводится сначала преобразованием простого произведения \(A^\mu B^\nu\). Как только правило известно для произведения векторов, линейность даёт правило для любого тензора второго ранга.

Для буста вдоль оси \(1\) координаты преобразуются так:
\begin{align}
  x'^1&=\gamma(x^1-vx^0), &
  x'^0&=\gamma(x^0-vx^1), \\
  x'^2&=x^2, & x'^3&=x^3.
\end{align}
Применим то же преобразование по отдельности к обоим индексам тензора. Для компоненты, параллельной бусту, на доске начинают с произведения:
\begin{align}
  (A'^0B'^1)
  =\gamma^2\bigl(
     A^0B^1-vA^0B^0-vA^1B^1+v^2A^1B^0
   \bigr).
\end{align}
Заменяя каждое произведение соответствующей компонентой тензора, а затем используя антисимметрию, \(F^{00}=F^{11}=0\) и \(F^{10}=-F^{01}\), получаем
\begin{equation}
  F'^{01}
  =\gamma^2(1-v^2)F^{01}
  =F^{01}.
\end{equation}
Следовательно, компонента электрического поля, параллельная движению, не изменяется. Параллельная компонента магнитного поля ведёт себя аналогично.

С поперечными компонентами дело обстоит иначе. Следуя конвенции штрихованной системы отсчёта и знаков, используемой на доске,
\begin{align}
  F'^{02}
  &=\gamma\bigl(F^{02}-vF^{12}\bigr), \\
  E'_2
  &=\gamma(E_2+vB_3).
  \label{eq:e-transform}
\end{align}
В поперечной компоненте электрического поля появляется вклад от магнитной компоненты. Остальные законы преобразования следуют из того же правила для двух индексов и оставлены в лекции как упражнение.

\begin{figure}[H]
\centering
\includegraphics[width=0.76\textwidth]{lecture_07_figure_13.png}
\caption{Преобразование поперечных компонент тензора в том месте, где видно, что компонента \(E'_2\) содержит член, пропорциональный \(vB_3\).}
\lectureframe{01:23:40}
\end{figure}

Знак в \eqref{eq:e-transform} зависит от того, какой системе отсчёта приписывают движение со скоростью \(+v\), а также от индексной и метрической конвенций. Но его инвариантное содержание от этого не меняется: поперечные электрические и магнитные компоненты смешиваются, а ковариантное уравнение силы остаётся справедливым после преобразования.

\begin{classroomqa}
\audiencequestion{Для тока в проводе можно ли двигаться с некоторой скоростью так, чтобы магнитное поле исчезло?}
\lecturerresponse{Недостаточно просто двигаться вместе с электронами проводимости, потому что положительные заряды решётки тогда движутся в новой системе отсчёта. Поэтому нейтральный провод с током — более тонкий случай, чем отдельный движущийся сгусток заряда. Чисто магнитное поле нельзя полностью устранить преобразованием. Если присутствуют и электрическое, и магнитное поля, то их преобразованные части при особых обстоятельствах могут взаимно компенсироваться; для одного равномерно движущегося заряда простой системой является система покоя заряда, в которой его магнитное поле исчезает.}
\end{classroomqa}

В точном инвариантном критерии используются
\begin{equation}
  \boldsymbol B^2-\boldsymbol E^2,
  \qquad
  \boldsymbol E\cdot\boldsymbol B.
\end{equation}
Система отсчёта, в которой \(\boldsymbol B'=0\), может существовать только для поля с преобладанием электрической составляющей и при \(\boldsymbol E\cdot\boldsymbol B=0\); поле, которое в некоторой системе отсчёта является чисто магнитным, имеет преобладание магнитной составляющей и не может удовлетворять этому условию.
\editorialnote{Эти инвариантные критерии точно выражают то уточнение, к которому пришли в ходе обсуждения в аудитории; их вывод откладывается.}

\begin{classroomqa}
\audiencequestion{Образуют ли электрический скалярный потенциал и магнитный векторный потенциал четырёхвектор?}
\lecturerresponse{Да. С точностью до знаков, зависящих от принятых соглашений, они объединяются в четырёхпотенциал \(A^\mu\). В уравнениях Максвелла этот объект будет использоваться систематически.}
\end{classroomqa}

\section{Производные полей}

Теперь мы оставим траекторию одной частицы и рассмотрим поле во всём пространстве-времени. Пусть
\begin{equation}
  \phi=\phi(x^1,x^2,x^3,x^0)
\end{equation}
будет скалярным полем. Его четыре производные образуют ковариантные компоненты четырёхвектора:
\begin{equation}
  \partial_\mu\phi
  \equiv\frac{\partial\phi}{\partial x^\mu}
  =\left(
    \frac{\partial\phi}{\partial x},
    \frac{\partial\phi}{\partial y},
    \frac{\partial\phi}{\partial z},
    \frac{\partial\phi}{\partial t}
  \right).
\end{equation}
Цепное правило, применённое к преобразованию координат Лоренца, доказывает этот ковариантный закон преобразования. Поднятие индекса меняет знаки у пространственных компонент на противоположные, но оставляет временной знак неизменным:
\begin{equation}
  \partial^\mu\phi
  =\left(
    -\frac{\partial\phi}{\partial x},
    -\frac{\partial\phi}{\partial y},
    -\frac{\partial\phi}{\partial z},
    \frac{\partial\phi}{\partial t}
  \right).
\end{equation}

Каждое дифференцирование добавляет нижний индекс. Поэтому
\(\partial_\nu V^\mu\) — тензор второго ранга с одним верхним и одним нижним индексом, а дифференцирование тензора даёт объект с одним дополнительным нижним индексом. Это основное правило, необходимое для записи уравнений поля без выбора выделенной инерциальной системы отсчёта.

\section{Инвариантное скалярное волновое уравнение}

Самые простые волновые уравнения содержат две производные. Если уравнение должно иметь один и тот же вид в каждой инерциальной системе отсчёта, его левая часть должна быть скаляром. Свёртка индексов производных даёт единственный простейший выбор:
\begin{equation}
  \boxed{\partial^\mu\partial_\mu\phi=0.}
  \label{eq:wave-covariant}
\end{equation}
Если расписать компоненты в используемой здесь сигнатуре,
\begin{equation}
  \frac{\partial^2\phi}{\partial t^2}
  -\frac{\partial^2\phi}{\partial x^2}
  -\frac{\partial^2\phi}{\partial y^2}
  -\frac{\partial^2\phi}{\partial z^2}=0.
  \label{eq:wave-components}
\end{equation}
Свёртка в \eqref{eq:wave-covariant} делает лоренц-инвариантность явной; уравнение \eqref{eq:wave-components} — это то же самое утверждение, записанное покомпонентно.

\begin{figure}[H]
\centering
\includegraphics[width=0.75\textwidth]{lecture_07_figure_14.png}
\caption{Форма скалярного волнового уравнения со свёрткой по индексам производных и её покомпонентное разложение вида «время минус пространство».}
\lectureframe{01:39:40}
\end{figure}

\begin{classroomqa}
\audiencequestion{Должна ли правая часть волнового уравнения обязательно быть нулевой?}
\lecturerresponse{Нет. Она может быть любым скаляром, совместимым с задачей. Константа, умноженная на \(\phi\), или нелинейная скалярная функция, например пропорциональная \(\phi^2\), также сохраняет ковариантность. Нуль даёт простейшее свободное волновое уравнение.}
\end{classroomqa}

Смысл этого скалярного примера — подготовка. В уравнениях Максвелла электрическое и магнитное поля смешиваются более сложным образом, но и их можно записать так, чтобы лоренц-инвариантность была видна из их индексной структуры. К этой задаче мы вернёмся в следующий раз.
